\subsection*{选择执行实例}

\co@inthreads@ 从当前并行作用域中选择执行实例。

下面只让索引不大于 2 的实例执行选定区域：

\begin{lstlisting}[language=croqtile,caption={tileflow 程序的并行视角}]
__co__  void foo() {
  // sequential code
  // start of the SPMD code region
  parallel p by 6 {
    // some uniform code for all threads
    // ...
    inthreads (p <= 2) {
      // divergent code for thread 0, 1, 2 (thread-mask: 000111)
    }
    // threads are forced to be synchronized
    // some uniform code for all threads
    // ...
  }
  // end of the SPMD code region
  // sequential code
}
\end{lstlisting}

可以把程序中的执行范围分成以下几类：

\begin{itemize}[leftmargin=1.6em]
\item 主机侧的顺序代码。
\item \co@parallel-by@ 内由多个实例执行的代码。
\item 并行实例共同执行的代码。
\item \co@inthreads@ 中仅由选定实例执行的代码。
\end{itemize}

本例的 \co@p <= 2@ 选择索引为 \co@0@、\co@1@、\co@2@ 的实例；其余实例跳过该区域。普通 \co@inthreads@ 用于按顺序执行的过滤区域，异步形式则使用 \co@.async@。

\co@inthreads@ 从已有并行区域中选择实例，使用位置和条件需要满足以下约束：

\begin{itemize}[leftmargin=1.6em]
\item 只能放在 \co@parallel-by@ 创建的并行作用域中。
\item 条件需要与并行索引有关。
\end{itemize}

编译器会检查条件表达式及其所在的并行作用域。
